\documentclass[../../main.tex]{subfiles}% 注意这里的文件路径不能用 ./main.tex ,否则用latexmk编译子文件会报错
\graphicspath{{\subfix{./image/}}{\subfix{../../image/}}} % 指定图片引用路径,后续可以直接使用图片文件名
% 如果图片存放在上级目录,则使用 ../../image/ 作为备用路径

% 例如:
% \begin{figure}[H]
% \centering
% \includegraphics[width=0.3\textwidth]{图.png}
% \caption{}
% \label{figure:图}
% \end{figure}
% 注意:上述\label{}一定要放在\caption{}之后,否则引用图片序号会只会显示??.

\begin{document}

\section{持续性}\label{section:ch06-持续性}

(P) 的极弱解一个值得注意的性质是其 $2$-集中部分（见\cref{definition:ch01-集中弥散}与\cref{subsection:ch03-集中分解}）在时间上持续：

\begin{theorem}\label{theorem:ch06-6.2.1}
设 (A0)--(A1) 成立，$u$ 为问题 (P) 的极弱解。则对每个 $t\in(0,T)$，有
\begin{align}
[u(\cdot,t)]_{c,2}=[u_0]_{c,2} \quad\text{在 }\mathscr{R}_f(\Omega). \label{align:ch6-6.12}
\end{align}
\end{theorem}

\begin{remark}\label{remark:ch06-6.2.1}
设 $u$ 为 (P) 的极弱解。由\cref{theorem:ch06-6.2.1}，对所有 $t\in(0,T)$，$u(\cdot,t)$ 是 $2$-弥散的当且仅当 $u_0$ 亦然。
\end{remark}

\begin{proof}[证明\cref{theorem:ch06-6.2.1}]
设 $K\subseteq\Omega$ 为任意紧集且 $C_2(K)=0$。由\cref{proposition:ch03-零容量逼近}存在序列 $\{f_n\}\subseteq C^\infty_c(\Omega)$，满足：(i) $f_n(\Omega)\subseteq[0,1]$，$f_n|_K=1$；(ii) $\|\Delta f_n\|_{L^1(\Omega)}\to0$（$n\to\infty$）。特别，因
\begin{align*}
\int_\Omega|\nabla f_n|^2\,\mathrm{d}x=-\int_\Omega f_n\Delta f_n\,\mathrm{d}x\leqslant\int_\Omega|\Delta f_n|\,\mathrm{d}x,
\end{align*}
有 $\lim_{n\to\infty}\|f_n\|_{H^1_0(\Omega)}=0$。

设 $U\subseteq\Omega$ 为开集且 $U\supseteq K$。由\cref{lemma:appA-光滑截断}与标准正则化论证，存在 $\rho_U\in C^\infty_c(U)$，$\rho_U(U)\subseteq[0,1]$，$\rho_U|_K=1$。对所有 $t\in(0,T]$ 及 $n\in\mathbb{N}$，由等式\eqref{align:ch6-6.4}得（见\eqref{align:ch03-3.82}）
\begin{align}
\begin{aligned}
&\langle[u(\cdot,t)]_{c,2},f_n\rho_U\rangle-\langle[u_0]_{c,2},f_n\rho_U\rangle\\
&= \int_0^t\mathrm{d}s\int_\Omega\phi(u_r(x,s))\Delta(f_n\rho_U)(x)\,\mathrm{d}x\\
&\quad+\langle[u_0]_{d,2},f_n\rho_U\rangle-\langle[u(\cdot,t)]_{d,2},f_n\rho_U\rangle.
\end{aligned} \label{align:ch6-6.13}
\end{align}
因 $\|f_n\|_{H^1_0(\Omega)}\to0$，有 $f_n\rho_U\stackrel{*}{\rightharpoonup}0$ 在 $L^\infty(\Omega)$，$f_n\rho_U\to0$ 在 $H^1_0(\Omega)$。又由\cref{theorem:ch03-扩散部分H负1}，$[u_0]_{d,2}$ 与 $[u(\cdot,t)]_{d,2}$（$t\in(0,T)$）均属于 $L^1(\Omega)+H^{-1}(\Omega)$。故
\begin{align}
\lim_{n\to\infty}\langle[u(\cdot,t)]_{d,2},f_n\rho_U\rangle
=\lim_{n\to\infty}\langle[u_0]_{d,2},f_n\rho_U\rangle=0
\quad(t\in(0,T)). \label{align:ch6-6.14}
\end{align}
另一方面，因 $\phi(u_r)\in L^\infty(Q)$，$\|\Delta f_n\|_{L^1(\Omega)}\to0$，$\|f_n\|_{H^1_0(\Omega)}\to0$，得
\begin{align}
\lim_{n\to\infty}\int_0^t\int_\Omega\phi(u_r)\Delta(f_n\rho_U)\,\mathrm{d}x\,\mathrm{d}s=0. \label{align:ch6-6.15}
\end{align}
注意因对所有 $n\in\mathbb{N}$ 有 $\chi_K\leqslant f_n\rho_U\leqslant\chi_U$，故对所有 $n\in\mathbb{N}$，
\begin{align}
[u_0]^\pm_{c,2}(K)\leqslant\langle[u_0]^\pm_{c,2},f_n\rho_U\rangle\leqslant[u_0]^\pm_{c,2}(U), \label{align:ch6-6.16a}\\
[u(\cdot,t)]^\pm_{c,2}(K)\leqslant\langle[u(\cdot,t)]^\pm_{c,2},f_n\rho_U\rangle\leqslant[u(\cdot,t)]^\pm_{c,2}(U)
\quad(t\in(0,T)). \label{align:ch6-6.16b}
\end{align}
又对任意 $t\in(0,T)$，
\begin{align}
\limsup_{n\to\infty}\langle[u(\cdot,t)]_{c,2},f_n\rho_U\rangle
\leqslant\limsup_{n\to\infty}\langle[u_0]_{c,2},f_n\rho_U\rangle, \label{align:ch6-6.17}
\end{align}
因为由\eqref{align:ch6-6.14}--\eqref{align:ch6-6.15}，\eqref{align:ch6-6.13}右端当 $n\to\infty$ 时趋于零。由\eqref{align:ch6-6.16a}--\eqref{align:ch6-6.17}，对任意 $t\in(0,T)$ 得
\begin{align}
\begin{aligned}
&[u(\cdot,t)]^+_{c,2}(K)-[u(\cdot,t)]^-_{c,2}(U)
\leqslant\langle[u(\cdot,t)]_{c,2},f_n\rho_U\rangle\\
&\quad\leqslant\limsup_{n\to\infty}\langle[u(\cdot,t)]_{c,2},f_n\rho_U\rangle
\leqslant\limsup_{n\to\infty}\langle[u_0]_{c,2},f_n\rho_U\rangle\\
&\quad\leqslant\limsup_{n\to\infty}\langle[u_0]^+_{c,2},f_n\rho_U\rangle
-\liminf_{n\to\infty}\langle[u_0]^-_{c,2},f_n\rho_U\rangle
\leqslant[u_0]^+_{c,2}(U)-[u_0]^-_{c,2}(K).
\end{aligned} \label{align:ch6-6.18}
\end{align}
类似地，
\begin{align}
[u_0]^+_{c,2}(K)-[u_0]^-_{c,2}(U)
\leqslant[u(\cdot,t)]^+_{c,2}(U)-[u(\cdot,t)]^-_{c,2}(K)
\quad(t\in(0,T)). \label{align:ch6-6.19}
\end{align}
由不等式\eqref{align:ch6-6.18}--\eqref{align:ch6-6.19}及测度 $[u_0]^\pm_{c,2}$、$[u(\cdot,t)]^\pm_{c,2}$ 的正则性（见\cref{lemma:ch01-正则性}），有
\begin{align*}
[u(\cdot,t)]^+_{c,2}(K)+[u_0]^-_{c,2}(K)
\leqslant[u_0]^+_{c,2}(K)+[u(\cdot,t)]^-_{c,2}(K)
\end{align*}
及
\begin{align*}
[u_0]^+_{c,2}(K)+[u(\cdot,t)]^-_{c,2}(K)
\leqslant[u(\cdot,t)]^+_{c,2}(K)+[u_0]^-_{c,2}(K).
\end{align*}
由这些不等式，对所有 $t\in(0,T)$，有
\begin{align*}
[u(\cdot,t)]_{c,2}(K)=[u_0]_{c,2}(K)
\quad\text{对每个 }2\text{-零容紧集 }K\subseteq\Omega,
\end{align*}
从而由 $K$ 的任意性及 $[u_0]_{c,2}$、$[u(\cdot,t)]_{c,2}$ 的正则性，
\begin{align}
[u(\cdot,t)]_{c,2}(E)=[u_0]_{c,2}(E)
\quad\text{对每个 }2\text{-零容集 }E\in\mathscr{B}(\Omega). \label{align:ch6-6.20}
\end{align}
另一方面，由 $2$-集中测度的定义，对每个 $t\in(0,T)$，存在 $2$-零容集 $E_1(t)\in\mathscr{B}(\Omega)$，使 $[u(\cdot,t)]_{c,2}=[u(\cdot,t)]_{c,2}\llcorner E_1(t)$。类似地，存在 $2$-零容集 $E_2\in\mathscr{B}(\Omega)$，使 $[u_0]_{c,2}=[u_0]_{c,2}\llcorner E_2$。由等式\eqref{align:ch6-6.20}，
\begin{align*}
[u(\cdot,t)]_{c,2}(E_1(t)\setminus E_2)=[u_0]_{c,2}(E_1(t)\setminus E_2)=0,
\end{align*}
\begin{align*}
[u_0]_{c,2}(E_2\setminus E_1(t))=[u(\cdot,t)]_{c,2}(E_2\setminus E_1(t))=0.
\end{align*}
因此 $[u(\cdot,t)]_{c,2}$ 与 $[u_0]_{c,2}$ 均集中于集 $F(t)\triangleq E_1(t)\cap E_2$。显然，对任意 $E\in\mathscr{B}(\Omega)$，有 $C_2(E\cap F(t))=C_2(F(t))=0$。由\eqref{align:ch6-6.20}，
\begin{align*}
[u(\cdot,t)]_{c,2}(E)=[u(\cdot,t)]_{c,2}(E\cap F(t))
=[u_0]_{c,2}(E\cap F(t))=[u_0]_{c,2}(E)
\end{align*}
对每个 $E\in\mathscr{B}(\Omega)$ 成立。结论随之成立。
\end{proof}




\end{document}